% phi3 2 punkt funktion
\documentclass{article}
\usepackage{/home/user1/code/tex/sty}
\usepackage{amsmath}
\usepackage{geometry}
\usepackage{hyperref}
\usepackage{graphicx}
\setlength{\parindent}{0pt}
\geometry{top=1in,bottom=1in,left=1in,right=1in}
\newcommand{\pic}[2]{\begin{center}\includegraphics[height=#1]{#2}\end{center}}
\begin{document}
Consider $\phi^3$ theory with mass.
\g{\mathcal{L}=\br{1}{2}\p_\mu \phi \p^\mu \phi - \br{1}{2}m^2\phi^2-\br{1}{3}g\phi^3}
The Feynman diagram is the 1-loop propagator with two external points.
\g{iM=\br{(-ig)^2}{2}\int \br{d^d k}{(2\pi)^d} \br{i}{k^2-m^2} \br{i}{(k+p)^2-m^2}}
Just evalute the integral part of the amplitude,
\g{\textrm{Integral}&=\int \br{d^d k}{(2\pi)^d} \br{i}{k^2-m^2} \br{i}{(k+p)^2-m^2}
\intertext{Insert Feynman parameters.}
&= \int_0^1 dx \int \br{d^dk}{(2\pi)^4} \br{1}{((1-x)(k^2-m^2)+x((k+p)^2-m^2) )^2}
\intertext{Simplify, and apply the momentum shift $l=k+xp$}
&= \int_0^1 dx  \int \br{d^dk}{(2\pi)^4} \br{1}{(l^2+x(1-x)p^2-m^2)^2} 
\intertext{Wick rotate to $l^0=il^0_E$ and apply the Euclidean spacetime integral.}
&= i \int_0^1 dx \br{\Gamma(2-\br{d}{2})}{(4\pi)^{d/2}} \br{1}{(m^2-x(1-x)p^2)^{2-d/2}}
\intertext{Let the denominator be $\Delta=m^2-x(1-x)p^2$.}
&= i \int_0^1 dx \br{\Gamma(2-\br{d}{2})}{(4\pi)^{d/2}} \br{1}{\Delta^{2-d/2}}
\intertext{Now apply $d=6-\epsilon$ everywhere.}
&= i \int_0^1 dx \br{\Gamma(\br{\epsilon}{2}-1)}{(4\pi)^{3-\br{\epsilon}{2}}} \br{1}{\Delta^{\br{\epsilon}{2}-1}}
\intertext{Expand the $\epsilon$ in all three terms, and keep only those proportional to $\br{1}{\epsilon}$ and a constant.}.
&=\br{i}{(4\pi)^3}\int_0^1 \Delta(-\br{2}{\epsilon}-1+\gamma-\log(4\pi)+\log(\Delta))
}
\end{document}
